% ============================================================
% FICHE SIMPLE — STS BIOAL 2e ANNÉE
% Les couples (moyenne ; écart type) et (médiane ; interquartile)
% Version vulgarisée + CORRECTION DÉTAILLÉE
% ============================================================

\documentclass[12pt,a4paper]{article}

\usepackage[utf8]{inputenc}
\usepackage[T1]{fontenc}
\usepackage[french]{babel}
\usepackage{amsmath,amssymb}
\usepackage{geometry}
\usepackage{enumitem}
\usepackage{booktabs}
\usepackage{array}
\usepackage{tabularx}
\usepackage{xcolor}
\usepackage{tcolorbox}
\usepackage{fancyhdr}
\usepackage{tikz}
\usepackage{pgfplots}
\pgfplotsset{compat=1.18}
\usetikzlibrary{babel}
\usetikzlibrary{arrows.meta}
\usepackage{lastpage}

\tcbuselibrary{skins,breakable}

\geometry{margin=2cm, headheight=15pt}

% ------------------------------------------------------------
% Couleurs
% ------------------------------------------------------------
\definecolor{primary}{RGB}{31,78,121}
\definecolor{soft}{RGB}{230,239,247}
\definecolor{amberbg}{RGB}{255,244,224}
\definecolor{amberborder}{RGB}{240,192,112}
\definecolor{greenbg}{RGB}{233,246,234}
\definecolor{greenborder}{RGB}{182,224,184}
\definecolor{redbg}{RGB}{253,236,236}
\definecolor{redborder}{RGB}{242,179,171}
\definecolor{purplebg}{RGB}{240,235,250}
\definecolor{purpleborder}{RGB}{190,170,220}
\definecolor{grayline}{RGB}{120,130,140}
\definecolor{accent}{RGB}{200,80,60}

% ------------------------------------------------------------
% En-tête / pied de page
% ------------------------------------------------------------
\pagestyle{fancy}
\fancyhf{}
\fancyhead[L]{\textcolor{primary}{\textbf{BTS BioALC — 2\textsuperscript{ème} année}}}
\fancyhead[R]{\textcolor{primary}{Fiche simple — Statistiques}}
\fancyfoot[L]{\textcolor{grayline}{\small Prénom : \dotfill}}
\fancyfoot[C]{\textcolor{grayline}{\small \thepage/\pageref{LastPage}}}
\fancyfoot[R]{\textcolor{grayline}{\small Fiche à garder}}
\renewcommand{\headrulewidth}{0.4pt}
\renewcommand{\footrulewidth}{0.4pt}

% ------------------------------------------------------------
% Boîtes
% ------------------------------------------------------------
\newtcolorbox{objectif}[1][]{
  colback=soft, colframe=primary, boxrule=0.8pt, arc=4pt,
  title={\textbf{But de cette fiche}}, fonttitle=\bfseries,
  coltitle=white, colbacktitle=primary,
  left=6pt, right=6pt, top=4pt, bottom=4pt,
  breakable, #1
}

\newtcolorbox{histoire}{
  colback=amberbg, colframe=amberborder, boxrule=0.8pt, arc=4pt,
  title={\textbf{L'histoire du jour}}, fonttitle=\bfseries,
  coltitle=black, colbacktitle=amberborder,
  left=6pt, right=6pt, top=4pt, bottom=4pt,
  breakable
}

\newtcolorbox{rappel}{
  colback=greenbg, colframe=greenborder, boxrule=0.8pt, arc=4pt,
  title={\textbf{À retenir}}, fonttitle=\bfseries,
  coltitle=black, colbacktitle=greenborder,
  left=6pt, right=6pt, top=4pt, bottom=4pt,
  breakable
}

\newtcolorbox{exemple}{
  colback=white, colframe=primary, boxrule=0.8pt, arc=4pt,
  title={\textbf{On calcule ensemble}}, fonttitle=\bfseries,
  coltitle=white, colbacktitle=primary,
  left=6pt, right=6pt, top=4pt, bottom=4pt,
  breakable
}

\newtcolorbox{enclair}{
  colback=purplebg, colframe=purpleborder, boxrule=0.8pt, arc=4pt,
  title={\textbf{En clair}}, fonttitle=\bfseries,
  coltitle=black, colbacktitle=purpleborder,
  left=6pt, right=6pt, top=4pt, bottom=4pt,
  breakable
}

\newtcolorbox{attention}{
  colback=redbg, colframe=redborder, boxrule=0.8pt, arc=4pt,
  title={\textbf{Attention}}, fonttitle=\bfseries,
  coltitle=black, colbacktitle=redborder,
  left=6pt, right=6pt, top=4pt, bottom=4pt,
  breakable
}

\newtcolorbox{graphbox}[1][]{
  colback=white, colframe=grayline, boxrule=0.6pt, arc=4pt,
  title={\textbf{Image à regarder}}, fonttitle=\bfseries,
  coltitle=white, colbacktitle=grayline,
  left=4pt, right=4pt, top=2pt, bottom=2pt,
  breakable, #1
}

\newtcolorbox{correctionbox}[1][]{
  colback=soft, colframe=primary, boxrule=1.2pt, arc=4pt,
  title={\textbf{CORRECTION}}, fonttitle=\bfseries,
  coltitle=white, colbacktitle=primary,
  left=6pt, right=6pt, top=4pt, bottom=4pt,
  breakable, #1
}

\newtcolorbox{rappelcible}[1][]{
  colback=greenbg, colframe=greenborder, boxrule=1pt, arc=4pt,
  title={\textbf{Rappel express — juste avant de calculer}}, fonttitle=\bfseries,
  coltitle=black, colbacktitle=greenborder,
  left=6pt, right=6pt, top=4pt, bottom=4pt,
  breakable, #1
}

% ------------------------------------------------------------
% Lignes pour écrire
% ------------------------------------------------------------
\newcommand{\lignes}[1]{%
  \vspace{0.15cm}%
  \foreach \x in {1,...,#1}{%
    \noindent\rule{\linewidth}{0.4pt}\par\vspace{0.6cm}%
  }%
}

\title{}
\author{}
\date{}

% ============================================================
\begin{document}
% ============================================================

% ------------------------------------------------------------
% Titre
% ------------------------------------------------------------
\begin{center}
  {\Large \textbf{Fiche simple de statistiques}}\\[0.3cm]
  {\large \textbf{La moyenne avec l'écart type}}\\[0.2cm]
  {\large \textbf{La médiane avec l'écart interquartile}}\\[0.4cm]
  {\large BTS BioALC — 2\textsuperscript{ème} année}
\end{center}

\vspace{0.4cm}

\begin{objectif}
Dans cette fiche, on va apprendre 4 mots :

\vspace{0.2cm}
\begin{enumerate}[leftmargin=*]
  \item La \textbf{moyenne}.
  \item L'\textbf{écart type}.
  \item La \textbf{médiane}.
  \item L'\textbf{écart interquartile}.
\end{enumerate}

\vspace{0.2cm}
On va aussi apprendre à \textbf{choisir} entre deux groupes :

\vspace{0.2cm}
\begin{itemize}[leftmargin=*]
  \item Groupe 1 : la moyenne et l'écart type.
  \item Groupe 2 : la médiane et l'écart interquartile.
\end{itemize}
\end{objectif}

\vspace{0.4cm}

\begin{histoire}
Vous travaillez dans un laboratoire. Vous mesurez le \textbf{glucose} dans le sang.

Vous faites \textbf{8 mesures} sur le même tube. Vous voulez savoir :

\vspace{0.2cm}
\begin{itemize}[leftmargin=*]
  \item Est-ce que mes mesures sont proches les unes des autres ?
  \item Ou est-ce qu'elles partent dans tous les sens ?
\end{itemize}

\vspace{0.2cm}
Pour répondre, on va utiliser la \textbf{moyenne} et l'\textbf{écart type}.
\end{histoire}

\vspace{0.5cm}

% ============================================================
\section*{Partie 1. La moyenne et l'écart type}
% ============================================================

\subsection*{1.1. C'est quoi, la moyenne ?}

\begin{rappel}
La \textbf{moyenne}, c'est le \textbf{centre} de la liste.

\vspace{0.2cm}
\textbf{Comment on fait ?}

\vspace{0.2cm}
\begin{enumerate}[leftmargin=*]
  \item On additionne toutes les valeurs.
  \item On divise par le nombre de valeurs.
\end{enumerate}
\end{rappel}

\begin{enclair}
\textbf{Exemple simple.} Vos notes : 10, 12, 14.

On additionne : $10 + 12 + 14 = 36$.

On divise par 3 (il y a 3 notes) : $36 \div 3 = 12$.

La moyenne est \textbf{12}.
\end{enclair}

\subsection*{1.2. C'est quoi, l'écart type ?}

\begin{rappel}
L'\textbf{écart type}, c'est un nombre qui dit \textbf{de combien les valeurs s'éloignent de la moyenne}.

\vspace{0.2cm}
\textbf{Idée simple :}
\begin{itemize}[leftmargin=*]
  \item Si l'écart type est \textbf{petit} : les valeurs sont \textbf{serrées} autour de la moyenne.
  \item Si l'écart type est \textbf{grand} : les valeurs sont \textbf{étalées}.
\end{itemize}
\end{rappel}

\begin{enclair}
\textbf{Comparez ces deux listes.}

\vspace{0.2cm}
Liste A : 9, 10, 11. Moyenne = 10. Les valeurs sont \textbf{serrées}. Écart type \textbf{petit}.

\vspace{0.2cm}
Liste B : 1, 10, 19. Moyenne = 10. Les valeurs sont \textbf{très étalées}. Écart type \textbf{grand}.

\vspace{0.2cm}
Les deux listes ont la même moyenne (10). Mais la liste B est beaucoup plus dispersée. C'est l'écart type qui le montre.
\end{enclair}

\subsection*{1.3. Comment on calcule l'écart type ?}

\begin{rappel}
On fait \textbf{4 étapes}. Pas plus.

\vspace{0.2cm}
\begin{enumerate}[leftmargin=*]
  \item On calcule la moyenne.
  \item Pour chaque valeur : on fait \textbf{valeur $-$ moyenne}. Puis on met le résultat \textbf{au carré} (on multiplie par lui-même).
  \item On additionne tous ces carrés. Puis on divise par (nombre de valeurs $- 1$).
  \item On prend la \textbf{racine carrée} du résultat.
\end{enumerate}

\vspace{0.2cm}
\textbf{Formule :}
\[
s = \sqrt{\frac{\text{somme des (valeur } - \text{ moyenne)}^2}{n - 1}}
\]
$n$ = nombre de valeurs.
\end{rappel}

\begin{attention}
\textbf{Ne paniquez pas devant la formule.}

\vspace{0.2cm}
On va la faire \textbf{ensemble}, étape par étape, dans l'exemple qui suit. Vous n'avez rien à mémoriser par cœur. Il faut juste comprendre les 4 étapes.
\end{attention}

\subsection*{1.4. Exemple pas à pas — 8 mesures de glucose}

\begin{exemple}
On a 8 mesures de glucose (en g/L) sur le même tube :

\begin{center}
\begin{tabular}{|c|c|c|c|c|c|c|c|c|}
\hline
Mesure & 1 & 2 & 3 & 4 & 5 & 6 & 7 & 8 \\
\hline
Valeur & 1,02 & 0,98 & 1,01 & 0,99 & 1,00 & 1,03 & 0,97 & 1,00 \\
\hline
\end{tabular}
\end{center}

\vspace{0.3cm}
\textbf{Étape 1 — On calcule la moyenne.}

On additionne tout :
\[
1{,}02 + 0{,}98 + 1{,}01 + 0{,}99 + 1{,}00 + 1{,}03 + 0{,}97 + 1{,}00 = 8{,}00
\]

On divise par 8 (il y a 8 mesures) :
\[
8{,}00 \div 8 = 1{,}00
\]

La moyenne est $1{,}00$ g/L.

\vspace{0.3cm}
\textbf{Étape 2 — Pour chaque valeur, on fait valeur $-$ moyenne. Puis on met au carré.}

La moyenne est 1,00. Donc on fait : valeur $- 1{,}00$. Puis on met au carré.

\begin{center}
\begin{tabular}{|c|c|c|}
\hline
Valeur & Valeur $-$ 1,00 & Résultat au carré \\
\hline
1,02 & $+0{,}02$ & $0{,}0004$ \\
\hline
0,98 & $-0{,}02$ & $0{,}0004$ \\
\hline
1,01 & $+0{,}01$ & $0{,}0001$ \\
\hline
0,99 & $-0{,}01$ & $0{,}0001$ \\
\hline
1,00 & $0{,}00$ & $0{,}0000$ \\
\hline
1,03 & $+0{,}03$ & $0{,}0009$ \\
\hline
0,97 & $-0{,}03$ & $0{,}0009$ \\
\hline
1,00 & $0{,}00$ & $0{,}0000$ \\
\hline
\end{tabular}
\end{center}

\vspace{0.3cm}
\textbf{Étape 3 — On additionne les carrés. Puis on divise par $8 - 1 = 7$.}

Somme des carrés :
\[
0{,}0004 + 0{,}0004 + 0{,}0001 + 0{,}0001 + 0 + 0{,}0009 + 0{,}0009 + 0 = 0{,}0028
\]

On divise par 7 :
\[
0{,}0028 \div 7 = 0{,}0004
\]

\vspace{0.3cm}
\textbf{Étape 4 — On prend la racine carrée.}

\[
\sqrt{0{,}0004} = 0{,}02
\]

\vspace{0.3cm}
\textbf{Résultat :}
\[
\boxed{\text{moyenne} = 1{,}00 \text{ g/L} \qquad \text{écart type} = 0{,}02 \text{ g/L}}
\]
\end{exemple}

\begin{enclair}
La moyenne est \textbf{1,00 g/L}. C'est le centre.

L'écart type est \textbf{0,02 g/L}. Cela veut dire :

\vspace{0.2cm}
\textbf{En moyenne, chaque mesure s'éloigne de 0,02 g/L par rapport à 1,00 g/L.}

\vspace{0.2cm}
Comme 0,02 est \textbf{petit}, les 8 mesures sont \textbf{très serrées}. L'appareil est \textbf{fiable}. C'est une bonne nouvelle pour le laboratoire.
\end{enclair}

\begin{graphbox}
\textbf{Image 1 — À quoi ressemble un petit écart type et un grand écart type.}

\vspace{0.3cm}
\begin{center}
\begin{tikzpicture}
\begin{axis}[
  width=12cm, height=5.5cm,
  axis lines=left,
  xlabel={\small Valeur},
  ylabel={\small Nombre de mesures},
  xtick=\empty,
  ytick=\empty,
  domain=-4:4,
  samples=120,
  smooth,
]
\addplot[primary, very thick] {exp(-x^2/0.5)};
\addplot[accent, very thick, dashed] {0.5*exp(-x^2/4)};
\node[primary, font=\small] at (axis cs:2.5,1.05) {écart type petit};
\node[accent, font=\small] at (axis cs:-2.5,0.55) {écart type grand};
\end{axis}
\end{tikzpicture}
\end{center}

\vspace{0.2cm}
{\small La courbe bleue (pleine) est \textbf{serrée} : petit écart type. La courbe rouge (pointillés) est \textbf{large} : grand écart type. Les deux ont le même centre.}
\end{graphbox}

% ============================================================
\newpage
% ============================================================

\section*{Partie 2. La médiane et l'écart interquartile}

\subsection*{2.1. C'est quoi, la médiane ?}

\begin{rappel}
La \textbf{médiane}, c'est la valeur qui \textbf{coupe la liste en deux}.

\vspace{0.2cm}
\textbf{Comment on fait ?}

\vspace{0.2cm}
\begin{enumerate}[leftmargin=*]
  \item On range les valeurs de la plus petite à la plus grande.
  \item On prend la valeur du milieu.
\end{enumerate}

\vspace{0.2cm}
Si la liste a un nombre \textbf{impair} de valeurs : la médiane est la valeur du milieu.

Si la liste a un nombre \textbf{pair} de valeurs : la médiane est la moyenne des deux valeurs du milieu.
\end{rappel}

\begin{enclair}
\textbf{Exemple 1 (impair).} Liste : 3, 5, 7, 9, 11.

Il y a 5 valeurs. La 3\textsuperscript{e} est au milieu. Donc médiane = \textbf{7}.

\vspace{0.2cm}
\textbf{Exemple 2 (pair).} Liste : 2, 4, 6, 8.

Il y a 4 valeurs. Les deux du milieu sont 4 et 6. On fait la moyenne : $(4 + 6) \div 2 = 5$. Donc médiane = \textbf{5}.
\end{enclair}

\subsection*{2.2. C'est quoi, les quartiles ?}

\begin{rappel}
On coupe la liste en \textbf{4 morceaux égaux}. On obtient 3 valeurs :

\vspace{0.2cm}
\begin{itemize}[leftmargin=*]
  \item $Q_1$ : le premier quartile. 25\,\% des valeurs sont en dessous.
  \item $Q_2$ : le deuxième quartile. C'est \textbf{la médiane}.
  \item $Q_3$ : le troisième quartile. 75\,\% des valeurs sont en dessous.
\end{itemize}
\end{rappel}

\begin{enclair}
Imaginez une file de 100 personnes rangées par taille.

\vspace{0.2cm}
\begin{itemize}[leftmargin=*]
  \item $Q_1$ = la taille de la 25\textsuperscript{e} personne.
  \item $Q_2$ (médiane) = la taille de la 50\textsuperscript{e} personne.
  \item $Q_3$ = la taille de la 75\textsuperscript{e} personne.
\end{itemize}

\vspace{0.2cm}
Entre $Q_1$ et $Q_3$, il y a \textbf{50 personnes sur 100}. C'est le \textbf{milieu de la file}.
\end{enclair}

\subsection*{2.3. C'est quoi, l'écart interquartile ?}

\begin{rappel}
L'\textbf{écart interquartile}, c'est une \textbf{soustraction} :

\[
\text{écart interquartile} = Q_3 - Q_1
\]

\vspace{0.2cm}
Il mesure la \textbf{largeur du milieu} de la liste.
\end{rappel}

\begin{enclair}
L'écart interquartile, c'est la \textbf{largeur du milieu} de la liste. On ignore les 25\,\% les plus petits et les 25\,\% les plus grands.

\vspace{0.2cm}
\textbf{Pourquoi c'est utile ?} Parce qu'on ignore les valeurs \textbf{bizarres} aux extrémités. Donc le résultat est plus \textbf{stable}.
\end{enclair}

\begin{graphbox}
\textbf{Image 2 — La boîte à moustaches : on voit tout d'un coup.}

\vspace{0.3cm}
\begin{center}
\begin{tikzpicture}[scale=1.0]
  \draw[thick, ->] (0,0) -- (12,0);
  \foreach \x/\lab in {1/2, 3/4, 5/6, 7/8, 9/10, 11/12} {
    \draw (\x,0) -- (\x,-0.15);
    \node[below, font=\small] at (\x,-0.15) {\lab};
  }
  \node[below, font=\small] at (6,-0.7) {Valeur};

  \draw[primary, very thick, fill=primary!15] (3,1.5) rectangle (9,2.5);
  \draw[primary, very thick] (3,1.2) -- (3,2.8);
  \draw[primary, very thick] (9,1.2) -- (9,2.8);
  \draw[accent, very thick] (6,1.5) -- (6,2.5);
  \draw[primary, thick] (1,2) -- (3,2);
  \draw[primary, thick] (9,2) -- (11,2);
  \draw[primary, thick] (1,1.7) -- (1,2.3);
  \draw[primary, thick] (11,1.7) -- (11,2.3);

  \node[above, primary, font=\small] at (3,2.8) {$Q_1$};
  \node[above, accent, font=\small] at (6,2.5) {médiane};
  \node[above, primary, font=\small] at (9,2.8) {$Q_3$};
  \node[below, primary, font=\small] at (1,1.7) {min};
  \node[below, primary, font=\small] at (11,1.7) {max};

  \draw[<->, thick, accent] (3,3.5) -- (9,3.5);
  \node[above, accent, font=\small] at (6,3.5) {écart interquartile = $Q_3 - Q_1$};
\end{tikzpicture}
\end{center}

\vspace{0.2cm}
{\small \textbf{Comment lire cette image ?} Le rectangle s'appelle la \textbf{boîte}. Sa largeur est l'\textbf{écart interquartile}. Le trait rouge au milieu est la \textbf{médiane}. Les traits à gauche et à droite sont les \textbf{moustaches}.}
\end{graphbox}

\subsection*{2.4. Exemple pas à pas — 11 boîtes de Petri}

\begin{exemple}
On mesure le diamètre de colonies (en mm) sur 11 boîtes.

\vspace{0.2cm}
\textbf{Étape 1 — On range les valeurs dans l'ordre.}

\begin{center}
\begin{tabular}{|c|c|c|c|c|c|c|c|c|c|c|c|}
\hline
N\textsuperscript{o} & 1 & 2 & 3 & 4 & 5 & 6 & 7 & 8 & 9 & 10 & 11 \\
\hline
Valeur & 3,8 & 4,1 & 4,5 & 4,9 & 5,2 & 5,3 & 5,6 & 6,0 & 6,4 & 7,1 & 9,8 \\
\hline
\end{tabular}
\end{center}

\vspace{0.3cm}
\textbf{Étape 2 — On trouve la médiane.}

Il y a 11 valeurs. C'est impair. La médiane est la 6\textsuperscript{e} valeur.

\[
\text{médiane} = 5{,}3 \text{ mm}
\]

\vspace{0.3cm}
\textbf{Étape 3 — On trouve $Q_1$.}

$Q_1$ = la 3\textsuperscript{e} valeur (25\,\% de 11 = 2,75, on arrondit à 3).

\[
Q_1 = 4{,}5 \text{ mm}
\]

\vspace{0.3cm}
\textbf{Étape 4 — On trouve $Q_3$.}

$Q_3$ = la 9\textsuperscript{e} valeur (75\,\% de 11 = 8,25, on arrondit à 9).

\[
Q_3 = 6{,}4 \text{ mm}
\]

\vspace{0.3cm}
\textbf{Étape 5 — On calcule l'écart interquartile.}

\[
6{,}4 - 4{,}5 = 1{,}9 \text{ mm}
\]

\vspace{0.3cm}
\textbf{Résultat :}
\[
\boxed{\text{médiane} = 5{,}3 \text{ mm} \qquad \text{écart interquartile} = 1{,}9 \text{ mm}}
\]
\end{exemple}

\begin{enclair}
La médiane est \textbf{5,3 mm}. Cela veut dire : la moitié des boîtes ont un diamètre plus petit que 5,3 mm, et l'autre moitié un diamètre plus grand.

\vspace{0.2cm}
L'écart interquartile est \textbf{1,9 mm}. Cela veut dire : le milieu de la liste (les 50\,\% du milieu) s'étale sur 1,9 mm.

\vspace{0.2cm}
\textbf{Remarquez :} il y a une valeur très bizarre, \textbf{9,8 mm}. Elle est très loin des autres. La médiane et l'écart interquartile \textbf{n'ont pas bougé} à cause d'elle. Ils décrivent le milieu de la liste. C'est ça, l'avantage.
\end{enclair}

\begin{graphbox}
\textbf{Image 3 — La boîte à moustaches de cet exemple.}

\vspace{0.3cm}
\begin{center}
\begin{tikzpicture}[scale=0.85]
  \draw[thick, ->] (0,0) -- (13,0);
  \foreach \x/\lab in {0.5/3, 2/4, 3.5/5, 5/6, 6.5/7, 8/8, 9.5/9, 11/10} {
    \draw (\x,0) -- (\x,-0.15);
    \node[below, font=\small] at (\x,-0.15) {\lab};
  }
  \node[below, font=\small] at (6.5,-0.7) {Diamètre (mm)};

  \draw[primary, very thick, fill=primary!15] (2.27,1.5) rectangle (4.36,2.5);
  \draw[primary, very thick] (2.27,1.2) -- (2.27,2.8);
  \draw[primary, very thick] (4.36,1.2) -- (4.36,2.8);
  \draw[accent, very thick] (3.15,1.5) -- (3.15,2.5);
  \draw[primary, thick] (1.5,2) -- (2.27,2);
  \draw[primary, thick] (4.36,2) -- (5.13,2);
  \draw[primary, thick] (1.5,1.7) -- (1.5,2.3);
  \draw[primary, thick] (5.13,1.7) -- (5.13,2.3);

  \draw[accent, very thick] (8.10,1.6) -- (8.10,2.4);
  \node[accent, above, font=\small] at (8.10,2.4) {9,8 = valeur bizarre};

  \node[above, primary, font=\small] at (2.27,2.8) {$Q_1$};
  \node[above, accent, font=\small] at (3.15,2.5) {médiane};
  \node[above, primary, font=\small] at (4.36,2.8) {$Q_3$};
\end{tikzpicture}
\end{center}

\vspace{0.2cm}
{\small La valeur 9,8 mm est \textbf{toute seule} à droite. On dit qu'elle est \textbf{bizarre} (ou aberrante). La médiane et l'écart interquartile ne sont pas perturbés par elle.}
\end{graphbox}

% ============================================================
\newpage
% ============================================================

\section*{Partie 3. Comment choisir entre les deux groupes ?}

\subsection*{3.1. Un tableau très simple}

\begin{center}
\renewcommand{\arraystretch}{1.6}
\begin{tabularx}{\linewidth}{|p{4cm}|X|X|}
\hline
& \textbf{Groupe 1} & \textbf{Groupe 2} \\
& \textbf{moyenne + écart type} & \textbf{médiane + écart interquartile} \\
\hline
Ce qu'on utilise & la moyenne et l'écart type & la médiane et l'écart interquartile \\
\hline
Si une valeur est bizarre & \textbf{le résultat change beaucoup} & \textbf{le résultat ne change presque pas} \\
\hline
Quand on l'utilise & quand les valeurs sont \textbf{régulières} & quand il y a des \textbf{valeurs bizarres} \\
\hline
Exemple au labo & vérifier qu'un appareil est fiable & décrire des patients très différents \\
\hline
\end{tabularx}
\end{center}

\subsection*{3.2. La règle à retenir}

\begin{rappel}
\textbf{Deux cas simples.}

\vspace{0.3cm}
\textbf{Cas 1 :} les valeurs sont \textbf{à peu près pareilles}. Pas de valeur bizarre.
\[
\Rightarrow \text{On utilise la moyenne et l'écart type.}
\]

\vspace{0.3cm}
\textbf{Cas 2 :} il y a une ou plusieurs valeurs \textbf{très différentes} des autres.
\[
\Rightarrow \text{On utilise la médiane et l'écart interquartile.}
\]

\vspace{0.3cm}
\textbf{Si on ne sait pas :} on calcule les deux. Puis on compare la moyenne et la médiane.

\vspace{0.2cm}
\begin{itemize}[leftmargin=*]
  \item Si la moyenne est \textbf{proche} de la médiane : les valeurs sont régulières.
  \item Si la moyenne est \textbf{très différente} de la médiane : il y a une valeur bizarre.
\end{itemize}
\end{rappel}

\begin{graphbox}
\textbf{Image 4 — Pourquoi la médiane ne bouge pas avec une valeur bizarre.}

\vspace{0.3cm}
\begin{center}
\begin{tikzpicture}[scale=0.85]
  \node[font=\small\bfseries, primary] at (5,4.2) {Série A — pas de valeur bizarre};
  \draw[thick, ->] (0,0) -- (11,0);
  \foreach \x in {3,4,4.5,5,5.5,6,6.5,7,7.5,8} {
    \fill[primary] (\x,1) circle (3pt);
  }
  \draw[accent, very thick] (5.5,0) -- (5.5,3);
  \node[accent, above, font=\small] at (5.5,3) {moyenne = médiane};

  \node[font=\small\bfseries, primary] at (5,-2) {Série B — avec une valeur bizarre à droite};
  \draw[thick, ->] (0,-6) -- (11,-6);
  \foreach \x in {3,4,4.5,5,5.5,6,6.5,7,7.5,8} {
    \fill[primary] (\x,-5) circle (3pt);
  }
  \fill[accent] (10.5,-5) circle (5pt);
  \node[accent, above, font=\small] at (10.5,-4.6) {valeur bizarre};

  \draw[accent, very thick] (6.5,-6) -- (6.5,-3);
  \node[accent, above, font=\small] at (6.5,-3) {moyenne décalée};

  \draw[primary, very thick] (5.7,-6) -- (5.7,-3);
  \node[primary, above, font=\small] at (5.7,-2.5) {médiane stable};
\end{tikzpicture}
\end{center}

\vspace{0.2cm}
{\small Dans la série A, la moyenne et la médiane sont au même endroit. Dans la série B, la valeur bizarre \textbf{tire la moyenne vers la droite}. Mais la médiane \textbf{ne bouge presque pas}.}
\end{graphbox}

% ============================================================
\newpage
% ============================================================

\section*{Partie 4. Un exemple complet avec 2 appareils}

\begin{exemple}
Le laboratoire compare 2 appareils pour mesurer le glucose (en g/L). On fait 10 mesures avec chacun.

\vspace{0.3cm}
\textbf{Appareil A :} 0,95 — 0,98 — 1,00 — 1,02 — 1,01 — 0,99 — 1,03 — 0,97 — 1,00 — 1,05

\vspace{0.2cm}
\textbf{Appareil B :} 0,80 — 0,90 — 0,95 — 0,98 — 1,00 — 1,02 — 1,05 — 1,10 — 1,20 — 1,30

\vspace{0.4cm}
\textbf{Résultats des calculs :}

\begin{center}
\renewcommand{\arraystretch}{1.5}
\begin{tabular}{|l|c|c|}
\hline
& \textbf{Appareil A} & \textbf{Appareil B} \\
\hline
Moyenne & 1,00 g/L & 1,03 g/L \\
\hline
Écart type & 0,03 g/L & 0,15 g/L \\
\hline
Médiane & 0,995 g/L & 1,01 g/L \\
\hline
Écart interquartile & 0,04 g/L & 0,15 g/L \\
\hline
\end{tabular}
\end{center}
\end{exemple}

\begin{enclair}
Regardez la colonne A et la colonne B.

\vspace{0.3cm}
\textbf{Appareil A} : écart type = 0,03 g/L. C'est \textbf{petit}. Les mesures sont \textbf{serrées}. L'appareil est \textbf{fiable}.

\vspace{0.3cm}
\textbf{Appareil B} : écart type = 0,15 g/L. C'est \textbf{5 fois plus grand}. Les mesures sont \textbf{étalées}. L'appareil est \textbf{moins fiable}.

\vspace{0.3cm}
\textbf{Conclusion :} on garde l'appareil A. L'appareil B doit être réparé ou réglé.
\end{enclair}

\begin{graphbox}
\textbf{Image 5 — Comparaison des deux appareils.}

\vspace{0.3cm}
\begin{center}
\begin{tikzpicture}[scale=0.9]
  \draw[thick, ->] (0,0) -- (13,0);
  \foreach \x/\lab in {0.5/0.80, 2/0.90, 3.5/1.00, 5/1.10, 6.5/1.20, 8/1.30} {
    \draw (\x,0) -- (\x,-0.15);
    \node[below, font=\small] at (\x,-0.15) {\lab};
  }
  \node[below, font=\small] at (6.5,-0.7) {Glucose (g/L)};

  \draw[primary, very thick, fill=primary!20] (4.5,2.5) rectangle (5.3,3.5);
  \draw[primary, very thick] (4.5,2.2) -- (4.5,3.8);
  \draw[primary, very thick] (5.3,2.2) -- (5.3,3.8);
  \draw[accent, very thick] (4.9,2.5) -- (4.9,3.5);
  \draw[primary, thick] (4.0,3) -- (4.5,3);
  \draw[primary, thick] (5.3,3) -- (6.0,3);
  \node[primary, font=\small\bfseries] at (1.5,3) {Appareil A};
  \node[primary, font=\small] at (5.0,4.3) {boîte étroite = fiable};

  \draw[primary, very thick, fill=primary!20] (4.0,0.5) rectangle (7.0,1.5);
  \draw[primary, very thick] (4.0,0.2) -- (4.0,1.8);
  \draw[primary, very thick] (7.0,0.2) -- (7.0,1.8);
  \draw[accent, very thick] (5.2,0.5) -- (5.2,1.5);
  \draw[primary, thick] (1.0,1) -- (4.0,1);
  \draw[primary, thick] (7.0,1) -- (11.0,1);
  \node[primary, font=\small\bfseries] at (1.5,1) {Appareil B};
  \node[primary, font=\small] at (6.5,2.3) {boîte large = pas fiable};
\end{tikzpicture}
\end{center}

\vspace{0.2cm}
{\small La boîte de l'appareil A est \textbf{étroite}. C'est bon signe. La boîte de l'appareil B est \textbf{large}. C'est mauvais signe.}
\end{graphbox}

% ============================================================
\newpage
% ============================================================

\section*{Partie 5. Ce qu'il faut retenir (très court)}

\begin{rappel}
\textbf{4 mots à connaître.}

\vspace{0.3cm}
\textbf{1. La moyenne.} C'est le centre de la liste. On additionne tout, on divise par le nombre de valeurs.

\vspace{0.3cm}
\textbf{2. L'écart type.} C'est de combien les valeurs s'éloignent de la moyenne, \textbf{en moyenne}. Petit écart type = valeurs serrées. Grand écart type = valeurs étalées.

\vspace{0.3cm}
\textbf{3. La médiane.} C'est la valeur du milieu. On range les valeurs, puis on prend celle du milieu.

\vspace{0.3cm}
\textbf{4. L'écart interquartile.} C'est $Q_3 - Q_1$. C'est la largeur du milieu de la liste. On ignore les valeurs extrêmes.

\vspace{0.4cm}
\textbf{2 groupes à connaître.}

\vspace{0.3cm}
\textbf{Groupe 1 :} moyenne + écart type. On l'utilise quand les valeurs sont \textbf{régulières}.

\vspace{0.3cm}
\textbf{Groupe 2 :} médiane + écart interquartile. On l'utilise quand il y a des \textbf{valeurs bizarres}.
\end{rappel}

\vspace{0.5cm}

\section*{Partie 6. Exercices pour s'entraîner}

\begin{exemple}
\textbf{Exercice 1.} Voici 9 valeurs :
\[
12{,}1 \quad 12{,}3 \quad 12{,}0 \quad 12{,}4 \quad 12{,}2 \quad 12{,}1 \quad 12{,}5 \quad 12{,}2 \quad 12{,}3
\]

\begin{enumerate}[leftmargin=*]
  \item Calculez la moyenne.
  \item Calculez la médiane.
  \item Regardez les valeurs : y a-t-il des valeurs bizarres ?
  \item Quel groupe utiliseriez-vous ? Justifiez.
\end{enumerate}

\lignes{8}
\end{exemple}

\vspace{0.4cm}

\begin{exemple}
\textbf{Exercice 2.} Voici 10 valeurs :
\[
28 \quad 31 \quad 29 \quad 30 \quad 32 \quad 45 \quad 30 \quad 29 \quad 31 \quad 30
\]

\begin{enumerate}[leftmargin=*]
  \item Rangez les valeurs dans l'ordre.
  \item Calculez la médiane.
  \item Regardez la valeur 45 : est-elle bizarre ?
  \item Quel groupe utiliseriez-vous pour décrire cette série ? Justifiez.
\end{enumerate}

\lignes{8}
\end{exemple}

\vspace{0.4cm}

\begin{exemple}
\textbf{Exercice 3.} Vrai ou faux ? Entourez la bonne réponse.

\vspace{0.3cm}
\begin{enumerate}[leftmargin=*]
  \item « L'écart type mesure la dispersion autour de la moyenne. » \hfill \textbf{VRAI} / \textbf{FAUX}
  \item « La médiane change beaucoup s'il y a une valeur bizarre. » \hfill \textbf{VRAI} / \textbf{FAUX}
  \item « L'écart interquartile est égal à $Q_3 - Q_1$. » \hfill \textbf{VRAI} / \textbf{FAUX}
  \item « On utilise la moyenne et l'écart type quand les valeurs sont régulières. » \hfill \textbf{VRAI} / \textbf{FAUX}
  \item « La médiane est toujours égale à la moyenne. » \hfill \textbf{VRAI} / \textbf{FAUX}
\end{enumerate}
\end{exemple}

% ============================================================
% ============================================================
%                  CORRECTION DÉTAILLÉE
% ============================================================
% ============================================================

\newpage

\begin{center}
  {\Large \textbf{CORRECTION DÉTAILLÉE}}\\[0.2cm]
  {\large \textbf{Exercices de la partie 6}}\\[0.2cm]
  {\large BTS BioALC — 2\textsuperscript{ème} année}
\end{center}

\vspace{0.3cm}

\begin{correctionbox}
\textbf{Comment utiliser cette correction ?}

\vspace{0.2cm}
\begin{itemize}[leftmargin=*]
  \item Lisez d'abord la question.
  \item Essayez de répondre seul(e) sur votre brouillon.
  \item Puis lisez la correction \textbf{étape par étape}.
  \item Les boîtes vertes « Rappel express » vous aident \textbf{au moment précis} où vous en avez besoin.
\end{itemize}
\end{correctionbox}

\vspace{0.5cm}

% ============================================================
\section*{Correction — Exercice 1}
% ============================================================

\textbf{Énoncé.} Voici 9 valeurs :
\[
12{,}1 \quad 12{,}3 \quad 12{,}0 \quad 12{,}4 \quad 12{,}2 \quad 12{,}1 \quad 12{,}5 \quad 12{,}2 \quad 12{,}3
\]

\vspace{0.3cm}

\begin{correctionbox}
\textbf{Question 1 — Calculez la moyenne.}
\end{correctionbox}

\begin{rappelcible}
\textbf{Rappel express — la moyenne.}

\vspace{0.2cm}
Pour calculer la moyenne :
\begin{enumerate}[leftmargin=*]
  \item On additionne \textbf{toutes} les valeurs.
  \item On divise par le \textbf{nombre} de valeurs.
\end{enumerate}

\vspace{0.2cm}
Formule : $\text{moyenne} = \dfrac{\text{somme des valeurs}}{\text{nombre de valeurs}}$.
\end{rappelcible}

\vspace{0.3cm}
\textbf{Étape 1 — On compte les valeurs.}

Il y a \textbf{9} valeurs. Donc $n = 9$.

\vspace{0.3cm}
\textbf{Étape 2 — On additionne toutes les valeurs.}

On y va doucement, une par une :

\begin{center}
\begin{tabular}{|c|c|c|}
\hline
\textbf{Addition} & \textbf{Résultat partiel} & \textbf{Explication} \\
\hline
$12{,}1$ & $12{,}1$ & on démarre \\
\hline
$+ 12{,}3$ & $24{,}4$ & \\
\hline
$+ 12{,}0$ & $36{,}4$ & \\
\hline
$+ 12{,}4$ & $48{,}8$ & \\
\hline
$+ 12{,}2$ & $61{,}0$ & \\
\hline
$+ 12{,}1$ & $73{,}1$ & \\
\hline
$+ 12{,}5$ & $85{,}6$ & \\
\hline
$+ 12{,}2$ & $97{,}8$ & \\
\hline
$+ 12{,}3$ & $\mathbf{110{,}1}$ & total \\
\hline
\end{tabular}
\end{center}

\textbf{Somme totale} $= 110{,}1$.

\vspace{0.3cm}
\textbf{Étape 3 — On divise par le nombre de valeurs.}

\[
\text{moyenne} = \frac{110{,}1}{9} \approx 12{,}23
\]

\vspace{0.3cm}
\textbf{Réponse :} la moyenne est $\boxed{12{,}23}$ (arrondie au centième).

\begin{enclair}
On peut vérifier que le résultat est \textbf{logique} : toutes les valeurs sont entre 12,0 et 12,5. La moyenne doit donc être entre 12,0 et 12,5. 12,23 est bien dans cet intervalle. C'est cohérent.
\end{enclair}

\vspace{0.5cm}

\begin{correctionbox}
\textbf{Question 2 — Calculez la médiane.}
\end{correctionbox}

\begin{rappelcible}
\textbf{Rappel express — la médiane.}

\vspace{0.2cm}
Pour calculer la médiane :
\begin{enumerate}[leftmargin=*]
  \item On \textbf{range} les valeurs de la plus petite à la plus grande.
  \item On prend la valeur \textbf{du milieu}.
\end{enumerate}

\vspace{0.2cm}
\textbf{Attention au nombre de valeurs :}
\begin{itemize}[leftmargin=*]
  \item Nombre \textbf{impair} : la médiane est la valeur du milieu (exactement).
  \item Nombre \textbf{pair} : la médiane est la moyenne des \textbf{deux} valeurs du milieu.
\end{itemize}
\end{rappelcible}

\vspace{0.3cm}
\textbf{Étape 1 — On range les valeurs dans l'ordre.}

On part de la liste de départ et on la réorganise :

\begin{center}
\begin{tabular}{|c|c|c|c|c|c|c|c|c|c|}
\hline
Position & 1 & 2 & 3 & 4 & 5 & 6 & 7 & 8 & 9 \\
\hline
Valeur rangée & 12,0 & 12,1 & 12,1 & 12,2 & \textbf{12,2} & 12,3 & 12,3 & 12,4 & 12,5 \\
\hline
\end{tabular}
\end{center}

\textbf{Vérification :} on a bien 9 valeurs. On n'en a oublié aucune.

\vspace{0.3cm}
\textbf{Étape 2 — On cherche la position du milieu.}

Il y a 9 valeurs. C'est \textbf{impair}. La valeur du milieu est à la position :
\[
\frac{9+1}{2} = \frac{10}{2} = 5
\]

Donc la médiane est la \textbf{5\textsuperscript{e} valeur}.

\vspace{0.3cm}
\textbf{Étape 3 — On lit la 5\textsuperscript{e} valeur.}

En regardant le tableau : la 5\textsuperscript{e} valeur est \textbf{12,2}.

\vspace{0.3cm}
\textbf{Réponse :} la médiane est $\boxed{12{,}2}$.

\begin{enclair}
La médiane est 12,2. Cela veut dire : il y a \textbf{4 valeurs} plus petites que 12,2, et \textbf{4 valeurs} plus grandes que 12,2. La valeur 12,2 est bien au milieu de la liste.
\end{enclair}

\vspace{0.5cm}

\begin{correctionbox}
\textbf{Question 3 — Regardez les valeurs : y a-t-il des valeurs bizarres ?}
\end{correctionbox}

\begin{rappelcible}
\textbf{Rappel express — valeur bizarre.}

\vspace{0.2cm}
Une \textbf{valeur bizarre} (on dit aussi \textbf{aberrante}), c'est une valeur \textbf{très éloignée} des autres.

\vspace{0.2cm}
\textbf{Comment la repérer ?} On regarde la liste rangée. Si une valeur est \textbf{toute seule} à un bout de la liste, très loin du groupe principal, c'est une valeur bizarre.
\end{rappelcible}

\vspace{0.3cm}
\textbf{Méthode — On regarde la liste rangée.}

\begin{center}
\begin{tabular}{|c|c|c|c|c|c|c|c|c|c|}
\hline
Position & 1 & 2 & 3 & 4 & 5 & 6 & 7 & 8 & 9 \\
\hline
Valeur rangée & 12,0 & 12,1 & 12,1 & 12,2 & 12,2 & 12,3 & 12,3 & 12,4 & 12,5 \\
\hline
\end{tabular}
\end{center}

\textbf{Analyse :}
\begin{itemize}[leftmargin=*]
  \item La plus petite valeur est 12,0.
  \item La plus grande valeur est 12,5.
  \item L'écart entre le minimum et le maximum est $12{,}5 - 12{,}0 = 0{,}5$.
\end{itemize}

Toutes les valeurs sont \textbf{très proches} les unes des autres. Aucune valeur n'est isolée à un bout de la liste.

\vspace{0.3cm}
\textbf{Réponse :} \textbf{non}, il n'y a pas de valeur bizarre. Toutes les valeurs sont comprises entre 12,0 et 12,5. C'est une série \textbf{régulière}.

\vspace{0.5cm}

\begin{correctionbox}
\textbf{Question 4 — Quel groupe utiliseriez-vous ? Justifiez.}
\end{correctionbox}

\begin{rappelcible}
\textbf{Rappel express — les deux groupes.}

\vspace{0.2cm}
\textbf{Groupe 1 :} moyenne + écart type. On l'utilise quand les valeurs sont \textbf{régulières} (pas de valeur bizarre).

\vspace{0.2cm}
\textbf{Groupe 2 :} médiane + écart interquartile. On l'utilise quand il y a des \textbf{valeurs bizarres}.

\vspace{0.2cm}
\textbf{Astuce :} si la moyenne et la médiane sont \textbf{proches}, la série est régulière. Si elles sont \textbf{très différentes}, il y a une valeur bizarre.
\end{rappelcible}

\vspace{0.3cm}
\textbf{Étape 1 — On compare la moyenne et la médiane.}

\begin{itemize}[leftmargin=*]
  \item Moyenne = 12,23.
  \item Médiane = 12,2.
\end{itemize}

Écart : $12{,}23 - 12{,}2 = 0{,}03$. C'est \textbf{très petit}.

\vspace{0.3cm}
\textbf{Étape 2 — On conclut.}

La moyenne et la médiane sont \textbf{quasiment identiques}. Cela confirme que la série est \textbf{régulière}, sans valeur bizarre.

\vspace{0.3cm}
\textbf{Réponse :} on utilise le \textbf{Groupe 1 : la moyenne et l'écart type}. C'est le groupe adapté quand les valeurs sont régulières et proches les unes des autres.

\vspace{0.3cm}
\textbf{Justification rédigée :}
\begin{quote}
\textit{« J'utilise la moyenne et l'écart type, car les 9 valeurs sont toutes comprises entre 12,0 et 12,5. Il n'y a pas de valeur bizarre. La moyenne (12,23) et la médiane (12,2) sont très proches, ce qui confirme que la série est régulière. »}
\end{quote}

% ============================================================
\newpage
% ============================================================

\section*{Correction — Exercice 2}

\textbf{Énoncé.} Voici 10 valeurs :
\[
28 \quad 31 \quad 29 \quad 30 \quad 32 \quad 45 \quad 30 \quad 29 \quad 31 \quad 30
\]

\vspace{0.3cm}

\begin{correctionbox}
\textbf{Question 1 — Rangez les valeurs dans l'ordre.}
\end{correctionbox}

\begin{rappelcible}
\textbf{Rappel express — ranger une liste.}

\vspace{0.2cm}
Ranger une liste, c'est mettre les valeurs de la plus \textbf{petite} à la plus \textbf{grande}.

\vspace{0.2cm}
\textbf{Astuce :} on peut compter combien de fois chaque valeur apparaît.

\vspace{0.2cm}
\textbf{Attention :} on ne supprime aucune valeur. On déplace juste les valeurs.
\end{rappelcible}

\vspace{0.3cm}
\textbf{Étape 1 — On repère les valeurs et leur nombre d'apparitions.}

\begin{center}
\begin{tabular}{|c|c|c|c|c|c|c|}
\hline
\textbf{Valeur} & 28 & 29 & 30 & 31 & 32 & 45 \\
\hline
\textbf{Nombre de fois} & 1 & 2 & 3 & 2 & 1 & 1 \\
\hline
\end{tabular}
\end{center}

\textbf{Vérification :} $1 + 2 + 3 + 2 + 1 + 1 = 10$. On a bien 10 valeurs.

\vspace{0.3cm}
\textbf{Étape 2 — On écrit la liste rangée.}

\begin{center}
\begin{tabular}{|c|c|c|c|c|c|c|c|c|c|c|}
\hline
Position & 1 & 2 & 3 & 4 & 5 & 6 & 7 & 8 & 9 & 10 \\
\hline
Valeur & 28 & 29 & 29 & 30 & 30 & 30 & 31 & 31 & 32 & 45 \\
\hline
\end{tabular}
\end{center}

\textbf{Réponse :} la liste rangée est : $28 \ ; \ 29 \ ; \ 29 \ ; \ 30 \ ; \ 30 \ ; \ 30 \ ; \ 31 \ ; \ 31 \ ; \ 32 \ ; \ 45$.

\vspace{0.5cm}

\begin{correctionbox}
\textbf{Question 2 — Calculez la médiane.}
\end{correctionbox}

\begin{rappelcible}
\textbf{Rappel express — médiane avec un nombre pair de valeurs.}

\vspace{0.2cm}
Quand il y a un nombre \textbf{pair} de valeurs, il n'y a pas \textbf{une seule} valeur au milieu : il y en a \textbf{deux}.

\vspace{0.2cm}
On fait alors la \textbf{moyenne des deux valeurs du milieu}.

\vspace{0.2cm}
Formule : $\text{médiane} = \dfrac{\text{valeur du milieu n\textsuperscript{o}1} + \text{valeur du milieu n\textsuperscript{o}2}}{2}$.
\end{rappelcible}

\vspace{0.3cm}
\textbf{Étape 1 — On repère les deux valeurs du milieu.}

Il y a 10 valeurs. C'est \textbf{pair}. Les deux valeurs du milieu sont à la \textbf{5\textsuperscript{e}} et à la \textbf{6\textsuperscript{e}} position.

\begin{center}
\begin{tabular}{|c|c|c|c|c|c|c|c|c|c|c|}
\hline
Position & 1 & 2 & 3 & 4 & \textbf{5} & \textbf{6} & 7 & 8 & 9 & 10 \\
\hline
Valeur & 28 & 29 & 29 & 30 & \textbf{30} & \textbf{30} & 31 & 31 & 32 & 45 \\
\hline
\end{tabular}
\end{center}

\vspace{0.3cm}
\textbf{Étape 2 — On fait la moyenne des deux valeurs du milieu.}

Les deux valeurs du milieu sont \textbf{30} et \textbf{30}.

\[
\text{médiane} = \frac{30 + 30}{2} = \frac{60}{2} = 30
\]

\vspace{0.3cm}
\textbf{Réponse :} la médiane est $\boxed{30}$.

\begin{enclair}
La médiane est 30. Cela veut dire : il y a \textbf{5 valeurs} inférieures ou égales à 30, et \textbf{5 valeurs} supérieures ou égales à 30. La valeur 30 coupe bien la liste en deux.
\end{enclair}

\vspace{0.5cm}

\begin{correctionbox}
\textbf{Question 3 — Regardez la valeur 45 : est-elle bizarre ?}
\end{correctionbox}

\begin{rappelcible}
\textbf{Rappel express — repérer une valeur bizarre.}

\vspace{0.2cm}
Une valeur est \textbf{bizarre} (aberrante) si elle est \textbf{très éloignée} du groupe principal.

\vspace{0.2cm}
\textbf{Méthode simple :} on regarde l'écart entre les valeurs normales et la valeur suspecte.

\vspace{0.2cm}
Si l'écart est \textbf{beaucoup plus grand} que les autres écarts, la valeur est bizarre.
\end{rappelcible}

\vspace{0.3cm}
\textbf{Étape 1 — On regarde les valeurs sans le 45.}

Sans le 45, les valeurs vont de 28 à 32. L'étendue est :
\[
32 - 28 = 4
\]

\vspace{0.3cm}
\textbf{Étape 2 — On regarde l'écart avec le 45.}

Le 45 est \textbf{13 unités} au-dessus de la plus grande valeur normale (32) :
\[
45 - 32 = 13
\]

\vspace{0.3cm}
\textbf{Étape 3 — On compare les deux écarts.}

\begin{itemize}[leftmargin=*]
  \item Étendue normale : 4.
  \item Écart avec le 45 : 13.
\end{itemize}

13 est \textbf{beaucoup plus grand} que 4. Le 45 est donc \textbf{très éloigné} des autres valeurs.

\vspace{0.3cm}
\textbf{Réponse :} \textbf{oui}, la valeur 45 est \textbf{bizarre} (aberrante).

\begin{enclair}
Le 45 est isolé tout seul à droite de la liste. Toutes les autres valeurs sont entre 28 et 32. Le 45 est donc une \textbf{anomalie}. Elle peut correspondre à une erreur de mesure, à un patient particulier, ou à un incident technique.
\end{enclair}

\vspace{0.5cm}

\begin{correctionbox}
\textbf{Question 4 — Quel groupe utiliseriez-vous pour décrire cette série ? Justifiez.}
\end{correctionbox}

\begin{rappelcible}
\textbf{Rappel express — choisir le bon groupe.}

\vspace{0.2cm}
\begin{itemize}[leftmargin=*]
  \item \textbf{Pas de valeur bizarre} $\rightarrow$ Groupe 1 (moyenne + écart type).
  \item \textbf{Il y a une valeur bizarre} $\rightarrow$ Groupe 2 (médiane + écart interquartile).
\end{itemize}

\vspace{0.2cm}
\textbf{Pourquoi ?} Parce que la moyenne et l'écart type sont \textbf{faussés} par les valeurs bizarres, alors que la médiane et l'écart interquartile \textbf{ne bougent presque pas}.
\end{rappelcible}

\vspace{0.3cm}
\textbf{Étape 1 — On rappelle ce qu'on a trouvé.}

À la question 3, on a montré que la valeur 45 est \textbf{bizarre}.

\vspace{0.3cm}
\textbf{Étape 2 — On en déduit le groupe à utiliser.}

Comme il y a une valeur bizarre, on doit utiliser le \textbf{Groupe 2 : la médiane et l'écart interquartile}.

\vspace{0.3cm}
\textbf{Pour information — Voyons ce qui se passerait avec la moyenne.}

Si on calculait la moyenne avec le 45 :
\[
\text{moyenne} = \frac{28+31+29+30+32+45+30+29+31+30}{10} = \frac{315}{10} = 31{,}5
\]

La moyenne (31,5) est \textbf{tirée vers le haut} par le 45. Elle ne représente pas bien le groupe principal (qui tourne autour de 30).

La \textbf{médiane} (30), elle, n'a pas bougé : elle décrit bien le cœur de la série.

\vspace{0.3cm}
\textbf{Réponse :} on utilise le \textbf{Groupe 2 : la médiane et l'écart interquartile}.

\vspace{0.3cm}
\textbf{Justification rédigée :}
\begin{quote}
\textit{« J'utilise la médiane et l'écart interquartile, car la série contient une valeur bizarre (45) très éloignée des autres. La moyenne serait faussée par cette valeur, alors que la médiane reste stable et décrit mieux le groupe principal. »}
\end{quote}

% ============================================================
\newpage
% ============================================================

\section*{Correction — Exercice 3}

\textbf{Énoncé.} Vrai ou faux ? Entourez la bonne réponse.

\vspace{0.4cm}

\begin{correctionbox}
\textbf{Question 1 — « L'écart type mesure la dispersion autour de la moyenne. »}
\end{correctionbox}

\begin{rappelcible}
\textbf{Rappel express — l'écart type.}

\vspace{0.2cm}
L'\textbf{écart type} est un nombre qui mesure \textbf{de combien les valeurs s'éloignent de la moyenne}, en moyenne.

\vspace{0.2cm}
\begin{itemize}[leftmargin=*]
  \item Écart type \textbf{petit} : les valeurs sont \textbf{serrées} autour de la moyenne.
  \item Écart type \textbf{grand} : les valeurs sont \textbf{étalées}.
\end{itemize}

\vspace{0.2cm}
Le mot « dispersion » veut dire « éloignement », « étalement ».
\end{rappelcible}

\vspace{0.3cm}
\textbf{Analyse :} la phrase dit exactement la bonne définition. L'écart type mesure bien la dispersion autour de la moyenne.

\vspace{0.3cm}
\textbf{Réponse :} \textbf{VRAI}.

\vspace{0.5cm}

\begin{correctionbox}
\textbf{Question 2 — « La médiane change beaucoup s'il y a une valeur bizarre. »}
\end{correctionbox}

\begin{rappelcible}
\textbf{Rappel express — la médiane est robuste.}

\vspace{0.2cm}
La médiane, c'est la valeur \textbf{du milieu}. Elle ne dépend que de la \textbf{position} des valeurs, pas de leur grandeur exacte.

\vspace{0.2cm}
Donc une valeur très grande ou très petite \textbf{ne déplace pas} la valeur du milieu.

\vspace{0.2cm}
\textbf{Exemple :} liste 1, 2, 3, 4, 100. La médiane est 3. Si on remplace 100 par 1000, la médiane reste 3.
\end{rappelcible}

\vspace{0.3cm}
\textbf{Analyse :} la phrase dit le \textbf{contraire} de la vérité. La médiane \textbf{ne change pas} (ou très peu) quand il y a une valeur bizarre.

\vspace{0.3cm}
\textbf{Réponse :} \textbf{FAUX}.

\vspace{0.5cm}

\begin{correctionbox}
\textbf{Question 3 — « L'écart interquartile est égal à $Q_3 - Q_1$. »}
\end{correctionbox}

\begin{rappelcible}
\textbf{Rappel express — l'écart interquartile.}

\vspace{0.2cm}
L'\textbf{écart interquartile}, c'est la \textbf{largeur du milieu} de la liste.

\vspace{0.2cm}
On le calcule avec une \textbf{soustraction} :
\[
\text{écart interquartile} = Q_3 - Q_1
\]

\vspace{0.2cm}
$Q_1$ = premier quartile (25\,\% en dessous). $Q_3$ = troisième quartile (75\,\% en dessous).
\end{rappelcible}

\vspace{0.3cm}
\textbf{Analyse :} la phrase donne exactement la bonne formule.

\vspace{0.3cm}
\textbf{Réponse :} \textbf{VRAI}.

\vspace{0.5cm}

\begin{correctionbox}
\textbf{Question 4 — « On utilise la moyenne et l'écart type quand les valeurs sont régulières. »}
\end{correctionbox}

\begin{rappelcible}
\textbf{Rappel express — quand utiliser le Groupe 1 ?}

\vspace{0.2cm}
On utilise la \textbf{moyenne et l'écart type} quand :
\begin{itemize}[leftmargin=*]
  \item les valeurs sont \textbf{régulières} (proches les unes des autres) ;
  \item il n'y a \textbf{pas de valeur bizarre}.
\end{itemize}
\end{rappelcible}

\vspace{0.3cm}
\textbf{Analyse :} la phrase dit exactement la bonne règle.

\vspace{0.3cm}
\textbf{Réponse :} \textbf{VRAI}.

\vspace{0.5cm}

\begin{correctionbox}
\textbf{Question 5 — « La médiane est toujours égale à la moyenne. »}
\end{correctionbox}

\begin{rappelcible}
\textbf{Rappel express — moyenne et médiane sont différentes.}

\vspace{0.2cm}
\begin{itemize}[leftmargin=*]
  \item La \textbf{moyenne} est le \textbf{centre calculé} (somme divisée par le nombre).
  \item La \textbf{médiane} est la \textbf{valeur du milieu} (après avoir rangé la liste).
\end{itemize}

\vspace{0.2cm}
Ces deux nombres sont \textbf{souvent proches}, mais \textbf{pas toujours égaux}.

\vspace{0.2cm}
\textbf{Exemple :} liste 1, 2, 3, 4, 100.
\begin{itemize}
  \item Moyenne = $(1+2+3+4+100) \div 5 = 110 \div 5 = 22$.
  \item Médiane = 3 (valeur du milieu).
\end{itemize}
22 et 3 sont très différents !
\end{rappelcible}

\vspace{0.3cm}
\textbf{Analyse :} la phrase utilise le mot « \textbf{toujours} ». C'est faux. La moyenne et la médiane ne sont égales que dans des cas particuliers (par exemple quand la liste est symétrique).

\vspace{0.3cm}
\textbf{Réponse :} \textbf{FAUX}.

\vspace{0.5cm}

% ============================================================
\section*{Récapitulatif des réponses}
% ============================================================

\begin{center}
\renewcommand{\arraystretch}{1.8}
\begin{tabular}{|c|c|}
\hline
\textbf{Question} & \textbf{Réponse} \\
\hline
Exercice 1 — Moyenne & 12,23 \\
\hline
Exercice 1 — Médiane & 12,2 \\
\hline
Exercice 1 — Valeurs bizarres ? & Non \\
\hline
Exercice 1 — Groupe choisi & Groupe 1 (moyenne + écart type) \\
\hline
Exercice 2 — Liste rangée & 28 ; 29 ; 29 ; 30 ; 30 ; 30 ; 31 ; 31 ; 32 ; 45 \\
\hline
Exercice 2 — Médiane & 30 \\
\hline
Exercice 2 — Le 45 est-il bizarre ? & Oui \\
\hline
Exercice 2 — Groupe choisi & Groupe 2 (médiane + écart interquartile) \\
\hline
Exercice 3 — Q1 & VRAI \\
\hline
Exercice 3 — Q2 & FAUX \\
\hline
Exercice 3 — Q3 & VRAI \\
\hline
Exercice 3 — Q4 & VRAI \\
\hline
Exercice 3 — Q5 & FAUX \\
\hline
\end{tabular}
\end{center}

\vspace{0.5cm}

\begin{enclair}
\textbf{Le message à retenir de ces corrections.}

\vspace{0.3cm}
\begin{enumerate}[leftmargin=*]
  \item On calcule toujours la \textbf{moyenne} en additionnant puis en divisant.
  \item On calcule toujours la \textbf{médiane} en rangeant puis en prenant la valeur du milieu.
  \item On repère une \textbf{valeur bizarre} en regardant si elle est très éloignée des autres.
  \item S'il n'y a \textbf{pas} de valeur bizarre $\rightarrow$ \textbf{Groupe 1}.
  \item S'il y a une valeur bizarre $\rightarrow$ \textbf{Groupe 2}.
\end{enumerate}
\end{enclair}

\vfill

\begin{center}
\textit{Correction détaillée — BTS BioALC — 2\textsuperscript{ème} année}\\
\textit{Document à conserver avec la fiche simple.}
\end{center}

\end{document}